% Options for packages loaded elsewhere
\PassOptionsToPackage{unicode}{hyperref}
\PassOptionsToPackage{hyphens}{url}
\PassOptionsToPackage{dvipsnames,svgnames,x11names}{xcolor}
%
\documentclass[
  letterpaper,
  DIV=11,
  numbers=noendperiod]{scrartcl}

\usepackage{amsmath,amssymb}
\usepackage{iftex}
\ifPDFTeX
  \usepackage[T1]{fontenc}
  \usepackage[utf8]{inputenc}
  \usepackage{textcomp} % provide euro and other symbols
\else % if luatex or xetex
  \usepackage{unicode-math}
  \defaultfontfeatures{Scale=MatchLowercase}
  \defaultfontfeatures[\rmfamily]{Ligatures=TeX,Scale=1}
\fi
\usepackage{lmodern}
\ifPDFTeX\else  
    % xetex/luatex font selection
\fi
% Use upquote if available, for straight quotes in verbatim environments
\IfFileExists{upquote.sty}{\usepackage{upquote}}{}
\IfFileExists{microtype.sty}{% use microtype if available
  \usepackage[]{microtype}
  \UseMicrotypeSet[protrusion]{basicmath} % disable protrusion for tt fonts
}{}
\makeatletter
\@ifundefined{KOMAClassName}{% if non-KOMA class
  \IfFileExists{parskip.sty}{%
    \usepackage{parskip}
  }{% else
    \setlength{\parindent}{0pt}
    \setlength{\parskip}{6pt plus 2pt minus 1pt}}
}{% if KOMA class
  \KOMAoptions{parskip=half}}
\makeatother
\usepackage{xcolor}
\usepackage[margin=1.0in,margin=1in]{geometry}
\setlength{\emergencystretch}{3em} % prevent overfull lines
\setcounter{secnumdepth}{5}
% Make \paragraph and \subparagraph free-standing
\makeatletter
\ifx\paragraph\undefined\else
  \let\oldparagraph\paragraph
  \renewcommand{\paragraph}{
    \@ifstar
      \xxxParagraphStar
      \xxxParagraphNoStar
  }
  \newcommand{\xxxParagraphStar}[1]{\oldparagraph*{#1}\mbox{}}
  \newcommand{\xxxParagraphNoStar}[1]{\oldparagraph{#1}\mbox{}}
\fi
\ifx\subparagraph\undefined\else
  \let\oldsubparagraph\subparagraph
  \renewcommand{\subparagraph}{
    \@ifstar
      \xxxSubParagraphStar
      \xxxSubParagraphNoStar
  }
  \newcommand{\xxxSubParagraphStar}[1]{\oldsubparagraph*{#1}\mbox{}}
  \newcommand{\xxxSubParagraphNoStar}[1]{\oldsubparagraph{#1}\mbox{}}
\fi
\makeatother


\providecommand{\tightlist}{%
  \setlength{\itemsep}{0pt}\setlength{\parskip}{0pt}}\usepackage{longtable,booktabs,array}
\usepackage{calc} % for calculating minipage widths
% Correct order of tables after \paragraph or \subparagraph
\usepackage{etoolbox}
\makeatletter
\patchcmd\longtable{\par}{\if@noskipsec\mbox{}\fi\par}{}{}
\makeatother
% Allow footnotes in longtable head/foot
\IfFileExists{footnotehyper.sty}{\usepackage{footnotehyper}}{\usepackage{footnote}}
\makesavenoteenv{longtable}
\usepackage{graphicx}
\makeatletter
\newsavebox\pandoc@box
\newcommand*\pandocbounded[1]{% scales image to fit in text height/width
  \sbox\pandoc@box{#1}%
  \Gscale@div\@tempa{\textheight}{\dimexpr\ht\pandoc@box+\dp\pandoc@box\relax}%
  \Gscale@div\@tempb{\linewidth}{\wd\pandoc@box}%
  \ifdim\@tempb\p@<\@tempa\p@\let\@tempa\@tempb\fi% select the smaller of both
  \ifdim\@tempa\p@<\p@\scalebox{\@tempa}{\usebox\pandoc@box}%
  \else\usebox{\pandoc@box}%
  \fi%
}
% Set default figure placement to htbp
\def\fps@figure{htbp}
\makeatother

\KOMAoption{captions}{tableheading}
\makeatletter
\@ifpackageloaded{tcolorbox}{}{\usepackage[skins,breakable]{tcolorbox}}
\@ifpackageloaded{fontawesome5}{}{\usepackage{fontawesome5}}
\definecolor{quarto-callout-color}{HTML}{909090}
\definecolor{quarto-callout-note-color}{HTML}{0758E5}
\definecolor{quarto-callout-important-color}{HTML}{CC1914}
\definecolor{quarto-callout-warning-color}{HTML}{EB9113}
\definecolor{quarto-callout-tip-color}{HTML}{00A047}
\definecolor{quarto-callout-caution-color}{HTML}{FC5300}
\definecolor{quarto-callout-color-frame}{HTML}{acacac}
\definecolor{quarto-callout-note-color-frame}{HTML}{4582ec}
\definecolor{quarto-callout-important-color-frame}{HTML}{d9534f}
\definecolor{quarto-callout-warning-color-frame}{HTML}{f0ad4e}
\definecolor{quarto-callout-tip-color-frame}{HTML}{02b875}
\definecolor{quarto-callout-caution-color-frame}{HTML}{fd7e14}
\makeatother
\makeatletter
\@ifpackageloaded{caption}{}{\usepackage{caption}}
\AtBeginDocument{%
\ifdefined\contentsname
  \renewcommand*\contentsname{Table of contents}
\else
  \newcommand\contentsname{Table of contents}
\fi
\ifdefined\listfigurename
  \renewcommand*\listfigurename{List of Figures}
\else
  \newcommand\listfigurename{List of Figures}
\fi
\ifdefined\listtablename
  \renewcommand*\listtablename{List of Tables}
\else
  \newcommand\listtablename{List of Tables}
\fi
\ifdefined\figurename
  \renewcommand*\figurename{Figure}
\else
  \newcommand\figurename{Figure}
\fi
\ifdefined\tablename
  \renewcommand*\tablename{Table}
\else
  \newcommand\tablename{Table}
\fi
}
\@ifpackageloaded{float}{}{\usepackage{float}}
\floatstyle{ruled}
\@ifundefined{c@chapter}{\newfloat{codelisting}{h}{lop}}{\newfloat{codelisting}{h}{lop}[chapter]}
\floatname{codelisting}{Listing}
\newcommand*\listoflistings{\listof{codelisting}{List of Listings}}
\makeatother
\makeatletter
\makeatother
\makeatletter
\@ifpackageloaded{caption}{}{\usepackage{caption}}
\@ifpackageloaded{subcaption}{}{\usepackage{subcaption}}
\makeatother

\usepackage{bookmark}

\IfFileExists{xurl.sty}{\usepackage{xurl}}{} % add URL line breaks if available
\urlstyle{same} % disable monospaced font for URLs
\hypersetup{
  pdftitle={Part 1a: Quantitative data collection},
  pdfauthor={Nathan Alexander, PhD},
  colorlinks=true,
  linkcolor={blue},
  filecolor={Maroon},
  citecolor={Blue},
  urlcolor={Blue},
  pdfcreator={LaTeX via pandoc}}


\title{Part 1a: Quantitative data collection}
\usepackage{etoolbox}
\makeatletter
\providecommand{\subtitle}[1]{% add subtitle to \maketitle
  \apptocmd{\@title}{\par {\large #1 \par}}{}{}
}
\makeatother
\subtitle{Quantitative Session of the 2026 HUSOE ELPS Methods and Data
Institute}
\author{Nathan Alexander, PhD}
\date{}

\begin{document}
\maketitle

\renewcommand*\contentsname{Table of contents}
{
\hypersetup{linkcolor=}
\setcounter{tocdepth}{3}
\tableofcontents
}

\section{What data collection
entails}\label{what-data-collection-entails}

Quantitative data collection is the systematic process of gathering
information that can be recorded, organized, counted, compared, and
analyzed.

Data collection is more than administering a survey or downloading a
spreadsheet. Before collecting data, researchers must make decisions
about:

\begin{itemize}
\tightlist
\item
  Who or what they want to study
\item
  Which questions they want to answer
\item
  How abstract ideas will be measured
\item
  Who is included, excluded, or difficult to reach
\item
  How information will be stored, protected, and shared
\item
  What conclusions the resulting data can reasonably support
\end{itemize}

\begin{center}\rule{0.5\linewidth}{0.5pt}\end{center}

A useful workflow is:

\[
\text{Research question} \rightarrow \text{Population} \rightarrow \text{Measures} \rightarrow \text{Sampling frame} \rightarrow \text{Sample} \rightarrow \text{Analysis}
\]

At every stage, researchers make choices and those choices affect what
the data represent and whose experiences may be absent from the
resulting analysis.

The handoff between theory and method also happens as these choices are
made.

\begin{center}\rule{0.5\linewidth}{0.5pt}\end{center}

\subsection{Activity: Start with a
question}\label{activity-start-with-a-question}

Imagine that an educational leadership program wants to understand
students' experiences of belonging in their schools.

Discuss:

\begin{enumerate}
\def\labelenumi{\arabic{enumi}.}
\tightlist
\item
  What is the population of interest?
\item
  What information would need to be collected?
\item
  Whose perspectives might be missing if the survey is distributed only
  by email?
\item
  What would count as an unethical use of the findings?
\end{enumerate}

\section{Quantification}\label{quantification}

\begin{itemize}
\item
  \textbf{Quantification} is the process of representing concepts,
  experiences, or events with numbers or categories that can be
  analyzed.
\item
  Quantification can make patterns visible, but it can also obscure
  context.

  -- For example, a suspension count may tell us how often suspensions
  occur, but it does not automatically explain why they occurred, how
  rules were enforced, or how students experienced the process.
\item
  Seminal articles on quantification have informed different
  perspectives of quantification, such as
  \href{https://www.jstor.org/stable/1049479}{Zuberi's ``Deracializing
  social statistics: Problems in the quantification of race''}.
\end{itemize}

\section{Quantification}\label{quantification-1}

For example, researchers may be interested in concepts such as:

\begin{itemize}
\tightlist
\item
  School belonging
\item
  Teacher support
\item
  Academic opportunity
\item
  Racial climate
\item
  Student discipline
\item
  Access to advanced coursework
\end{itemize}

These concepts are not directly observed in the same way as height or
age.

\section{Quantification}\label{quantification-2}

Researchers must decide how to \textbf{operationalize}: that is,
translate a concept into observable indicators.

\subsection{Questions to ask before
measuring}\label{questions-to-ask-before-measuring}

\begin{itemize}
\tightlist
\item
  Who defines the concept?
\item
  What exactly does this variable represent?
\item
  Who defined the categories?
\item
  Do participants understand the question in the same way?
\item
  Does the measure capture the full experience we care about?
\item
  Could the measure reinforce stereotypes or harm a community?
\item
  What contextual information is needed to interpret the numbers
  responsibly?
\end{itemize}

\section{Sample concept to survey
quantifications}\label{sample-concept-to-survey-quantifications}

\begin{longtable}[]{@{}
  >{\raggedright\arraybackslash}p{(\linewidth - 4\tabcolsep) * \real{0.3333}}
  >{\raggedright\arraybackslash}p{(\linewidth - 4\tabcolsep) * \real{0.3333}}
  >{\raggedright\arraybackslash}p{(\linewidth - 4\tabcolsep) * \real{0.3333}}@{}}
\toprule\noalign{}
\begin{minipage}[b]{\linewidth}\raggedright
Concept
\end{minipage} & \begin{minipage}[b]{\linewidth}\raggedright
Possible measure
\end{minipage} & \begin{minipage}[b]{\linewidth}\raggedright
Example survey item
\end{minipage} \\
\midrule\noalign{}
\endhead
\bottomrule\noalign{}
\endlastfoot
School belonging & Agreement scale & ``I feel like I am part of this
school community.'' \\
Teacher support & Frequency or agreement scale & ``My teachers believe I
can be successful.'' \\
Access to opportunity & Course enrollment or availability & ``Is
Advanced Placement Calculus offered at your school?'' \\
Racial climate & Agreement scale & ``Students from different racial and
ethnic backgrounds are treated fairly at this school.'' \\
Exclusionary discipline & Administrative-record count & Number of
suspensions during an academic year \\
\end{longtable}

\begin{center}\rule{0.5\linewidth}{0.5pt}\end{center}

\subsection{Types of studies and sampling
strategies}\label{types-of-studies-and-sampling-strategies}

The methods used to collect sample data for statistical analysis is
extremely important.

If sample data are not collected appropriately, resulting statistical
analyses will be futile.

As a result, planning a study by identifying \emph{research questions},
the \emph{population} and \emph{sample} of interest, and selecting the
appropriate \emph{research method(s)} that will be used to \emph{analyze
data} that is collected are all essential parts in the statistical data
analysis process.

\begin{center}\rule{0.5\linewidth}{0.5pt}\end{center}

\subsubsection{Understanding experimental and observational study
designs}\label{understanding-experimental-and-observational-study-designs}

There are many different types of research studies.

Some studies use non-traditional methods (such as oral traditions) to
collect data, while others focus on more traditional methods (such as
surveys) to analyze data on a sample or a population.

These data collection methods produce a set of observations upon which
statistical analyses can be applied. We consider two core study designs
in statistical data analysis: \emph{experimental} studies and
\emph{observational} studies.

\begin{center}\rule{0.5\linewidth}{0.5pt}\end{center}

\textbf{Experimental study}: In an experimental study, a
\emph{treatment} is applied to a sample of interest to observe its
effects. There is generally a control group and a treatment group used
to understand the effects of the treatment. Individual observations are
referred to as experimental units whereas studies involving humans are
generally defined as study subjects.

\textbf{Observational study}: In an observational study, specific
characteristics of a sample or population are observed and measured but
individual observations or subjects of study are not influenced or
modified in any way.

\begin{center}\rule{0.5\linewidth}{0.5pt}\end{center}

\paragraph{Three types of observational
studies}\label{three-types-of-observational-studies}

\begin{tcolorbox}[enhanced jigsaw, opacityback=0, opacitybacktitle=0.6, coltitle=black, title={\textbf{DEFINITIONS}: Types of studies}, colback=white, arc=.35mm, breakable, bottomtitle=1mm, colbacktitle=quarto-callout-caution-color!10!white, toptitle=1mm, colframe=quarto-callout-caution-color-frame, titlerule=0mm, rightrule=.15mm, bottomrule=.15mm, toprule=.15mm, leftrule=.75mm, left=2mm]

\textbf{Retrospective study}: In a retrospective study, we go back in
time to collect data over some past period.

\textbf{Cross-sectional study}: In a cross-sectional study, data are
collected and measured at one point in time.

\textbf{Prospective study}: In a prospective study, we set up a study to
go forward in time and observe groups sharing common factors.

\end{tcolorbox}

\begin{center}\rule{0.5\linewidth}{0.5pt}\end{center}

\paragraph{Identify and describe the different sampling
methods}\label{identify-and-describe-the-different-sampling-methods}

There are two broad categories of selecting members of a population to
generate sample data:

\begin{itemize}
\item
  Probability sampling
\item
  Non-probability sampling
\end{itemize}

Within these two broad categories are other methods based on the needs
of the study. Each methods is used to support statistical data analysis
with some methods providing stronger evidence than others.

\section{Sampling data}\label{sampling-data}

Most researchers cannot collect data from every member of a population.
Instead, they collect data from a \textbf{sample} and use that
information to describe or learn about a broader population. This is
non-probability sampling.

\begin{itemize}
\tightlist
\item
  \textbf{Population:} The complete group a researcher wants to
  understand
\item
  \textbf{Sampling frame:} The list, system, or source from which the
  sample is selected
\item
  \textbf{Sample:} The individuals, schools, districts, documents, or
  other units actually included in the study
\item
  \textbf{Respondents:} Sampled people who ultimately provide data
\end{itemize}

\section{Sampling data}\label{sampling-data-1}

For example:

\begin{longtable}[]{@{}
  >{\raggedright\arraybackslash}p{(\linewidth - 2\tabcolsep) * \real{0.5000}}
  >{\raggedright\arraybackslash}p{(\linewidth - 2\tabcolsep) * \real{0.5000}}@{}}
\toprule\noalign{}
\begin{minipage}[b]{\linewidth}\raggedright
Element
\end{minipage} & \begin{minipage}[b]{\linewidth}\raggedright
Example
\end{minipage} \\
\midrule\noalign{}
\endhead
\bottomrule\noalign{}
\endlastfoot
Research question & How do high school students experience academic
advising? \\
Population & All high school students at a single school \\
Sampling frame & The high school student enrollment lists \\
Sample & 500 students randomly selected from that list \\
Respondents & The students who complete the survey \\
\end{longtable}

\begin{center}\rule{0.5\linewidth}{0.5pt}\end{center}

\subsection{The sample frame}\label{the-sample-frame}

The \textbf{sampling frame} is the practical source used to identify and
contact members of the target population. Examples include:

\begin{itemize}
\tightlist
\item
  A student enrollment roster
\item
  A list of school employees
\item
  A voter-registration list
\item
  A set of residential addresses
\item
  An online research panel
\item
  A list of organizations or institutions
\item
  Publicly available social-media posts matching a defined criterion
\end{itemize}

A frame is not automatically the same as the population. For example, an
institutional email list may exclude students with outdated contact
information, students who have disengaged from the institution, or
individuals who have limited internet access.

\begin{center}\rule{0.5\linewidth}{0.5pt}\end{center}

\subsubsection{Comprehensiveness}\label{comprehensiveness}

A strong sampling frame covers as much of the target population as
possible.

Two common concerns are:

\begin{itemize}
\tightlist
\item
  \textbf{Undercoverage:} Some members of the target population are
  absent from the frame
\item
  \textbf{Overcoverage:} The frame includes people who are not part of
  the target population
\end{itemize}

For example, a survey of current students using an old enrollment list
may include graduates while missing newly enrolled students.

\textbf{Prompt:} If you want to study student experiences with remote
learning, who might be missing from an email-based survey frame?

\begin{center}\rule{0.5\linewidth}{0.5pt}\end{center}

\subsubsection{Probability of selection}\label{probability-of-selection}

In a \textbf{probability sample}, each eligible unit has a known or
estimable, nonzero chance of selection.

Common probability-sampling approaches include:

\begin{itemize}
\tightlist
\item
  \textbf{Simple random sampling:} Every unit has the same probability
  of selection
\item
  \textbf{Systematic sampling:} Select every (k)-th unit after a random
  starting point
\item
  \textbf{Stratified sampling:} Divide the population into groups, then
  sample within each group
\item
  \textbf{Cluster sampling:} Sample groups, such as schools or
  classrooms, rather than individuals first
\item
  \textbf{Multistage sampling:} Select units in stages, such as
  districts, then schools, then students
\end{itemize}

Probability sampling supports statistical inference because the
selection process is documented. It does not guarantee a perfect sample;
nonresponse, measurement error, and incomplete frames can still
introduce bias.

\begin{center}\rule{0.5\linewidth}{0.5pt}\end{center}

\subsubsection{Efficiency}\label{efficiency}

Sampling design involves tradeoffs among cost, time, precision,
feasibility, and fairness.

For example, it may be efficient to survey one classroom, but that
classroom may not represent all students in a school. A stratified
sample may require more planning, but it can ensure that smaller groups
are adequately represented.

Researchers should ask:

\begin{itemize}
\tightlist
\item
  Which groups must be represented in the study?
\item
  Which groups may be difficult to reach?
\item
  What resources are available?
\item
  What degree of precision is necessary?
\item
  What harms or burdens could data collection impose?
\end{itemize}

\begin{center}\rule{0.5\linewidth}{0.5pt}\end{center}

\subsection{Example: Traditionalist
view}\label{example-traditionalist-view}

Suppose a school district asks:

\begin{quote}
What percentage of students report that they feel safe at school?
\end{quote}

A traditionalist quantitative approach might:

\begin{itemize}
\tightlist
\item
  Define the population as all enrolled students
\item
  Use the enrollment list as the sampling frame
\item
  Draw a random sample of students
\item
  Administer a standardized survey
\item
  Estimate the percentage of students who report feeling safe
\item
  Report a margin of error or confidence interval
\end{itemize}

This approach prioritizes consistency, comparability, and generalization
from the sample to the population.

\begin{center}\rule{0.5\linewidth}{0.5pt}\end{center}

\subsection{Example: Deficit view}\label{example-deficit-view}

A deficit-oriented approach might ask:

\begin{quote}
Why do students from a particular neighborhood report lower feelings of
safety?
\end{quote}

The problem is not the use of numbers. The concern is the explanation
implied by the question. If researchers immediately attribute
differences to students, families, cultures, or neighborhoods, they can
overlook unequal school resources, discriminatory discipline,
neighborhood disinvestment, violence exposure, or institutional policy.

A more careful researcher would ask:

\begin{itemize}
\tightlist
\item
  What conditions shape students' experiences of safety?
\item
  How are school rules and disciplinary practices implemented?
\item
  Are all groups equally protected and heard?
\item
  What institutional resources are available across schools?
\end{itemize}

\begin{center}\rule{0.5\linewidth}{0.5pt}\end{center}

\subsection{Sampling procedures}\label{sampling-procedures}

Researchers select samples using procedures that fit their research
question, population, and practical constraints.

\textbf{Quick question}: \emph{What are some sampling procedures that
you are familiar with or will be using for your study?}

\begin{center}\rule{0.5\linewidth}{0.5pt}\end{center}

\subsection{Sampling procedures}\label{sampling-procedures-1}

Researchers select samples using procedures that fit their research
question, population, and practical constraints.

\textbf{Quick question}: \emph{What are some sampling procedures that
you are familiar with or will be using for your study?}

Convenience and purposive samples can be useful, especially for
exploratory work, program feedback, and hard-to-reach populations.
However, researchers should avoid claiming that such samples
statistically represent a larger population unless the design supports
that conclusion.

\begin{center}\rule{0.5\linewidth}{0.5pt}\end{center}

\subsection{Sampling procedures}\label{sampling-procedures-2}

\begin{longtable}[]{@{}
  >{\raggedright\arraybackslash}p{(\linewidth - 4\tabcolsep) * \real{0.3333}}
  >{\raggedright\arraybackslash}p{(\linewidth - 4\tabcolsep) * \real{0.3333}}
  >{\raggedright\arraybackslash}p{(\linewidth - 4\tabcolsep) * \real{0.3333}}@{}}
\toprule\noalign{}
\begin{minipage}[b]{\linewidth}\raggedright
Procedure
\end{minipage} & \begin{minipage}[b]{\linewidth}\raggedright
Description
\end{minipage} & \begin{minipage}[b]{\linewidth}\raggedright
Example use
\end{minipage} \\
\midrule\noalign{}
\endhead
\bottomrule\noalign{}
\endlastfoot
Convenience sampling & Recruit people who are easiest to access & An
optional workshop exit survey \\
Simple random sampling & Randomly select individuals from a complete
list & Select 200 students from an enrollment roster \\
Stratified sampling & Sample separately within key groups & Ensure
representation across class year or school level \\
Cluster sampling & Sample groups, then individuals within groups &
Randomly select schools, then classrooms \\
Purposive sampling & Select cases meeting specified criteria & Interview
principals serving rural districts \\
Snowball sampling & Participants recruit additional participants & Reach
a hard-to-contact community or professional network \\
\end{longtable}

\begin{center}\rule{0.5\linewidth}{0.5pt}\end{center}

\subsection{Example: Traditionalist
view}\label{example-traditionalist-view-1}

A researcher wants to estimate the proportion of doctoral students who
use campus mental-health services. A probability-oriented design could:

\begin{enumerate}
\def\labelenumi{\arabic{enumi}.}
\tightlist
\item
  Obtain a list of currently enrolled doctoral students
\item
  Stratify the list by program or enrollment status
\item
  Randomly select students within each group
\item
  Invite selected students to complete a confidential survey
\item
  Calculate weighted estimates if groups were sampled at different rates
\item
  Clearly report the response rate, sampling procedure, and limitations
\end{enumerate}

\begin{center}\rule{0.5\linewidth}{0.5pt}\end{center}

\subsection{Example: Critical
quantification}\label{example-critical-quantification}

A critical approach begins by asking not only, ``How can we sample
efficiently?'' but also:

\begin{itemize}
\tightlist
\item
  Who has power to define the research problem?
\item
  Who benefits from the data collection?
\item
  Whose knowledge is treated as credible?
\item
  Which groups are likely to be excluded by the sampling frame?
\item
  How will findings be returned to participants and communities?
\item
  Could the results be used to stigmatize, surveil, or penalize people?
\end{itemize}

For example, a study of student belonging could involve students in
developing survey items, include accessible modes of participation
beyond email, report results back to students, and connect findings to
institutional action. This does not mean abandoning rigor. It means
treating validity, representation, transparency, and accountability as
connected concerns.

\begin{center}\rule{0.5\linewidth}{0.5pt}\end{center}

\subsection{Sample design and
estimates}\label{sample-design-and-estimates}

A sample design is the plan for selecting units from the population.
That design determines what estimates can be made and how confidently
they can be interpreted.

For a population quantity such as the proportion of students who feel
safe at school, we may use a sample estimate:

\[
\hat{p} = \frac{\text{Number of sampled respondents reporting safety}}{\text{Number of sampled respondents}}
\]

The estimate \(\hat{p}\) is not automatically the true population value.
It is an estimate based on the respondents and the design used to select
them.

\begin{center}\rule{0.5\linewidth}{0.5pt}\end{center}

\subsection{Sample design and
estimates}\label{sample-design-and-estimates-1}

When interpreting estimates, consider:

\begin{itemize}
\tightlist
\item
  Sampling error: Results would vary somewhat if a different sample were
  selected
\item
  Nonresponse: Selected people may not participate
\item
  Coverage error: Some people may be absent from the sampling frame
\item
  Measurement error: Questions may be unclear, sensitive, or interpreted
  differently
\item
  Weighting: Some samples require weights to reflect unequal
  probabilities of selection or differential response
\end{itemize}

\textbf{Key idea:} A larger sample is helpful, but sample quality
depends on more than sample size. A very large convenience sample can
still be biased.

\section{Storing data}\label{storing-data}

Data storage is part of research design, not an afterthought.
Researchers have a responsibility to protect participants, document
their work, and make appropriate decisions about access and sharing.
Good data-management practices include:

\begin{itemize}
\tightlist
\item
  Store identifiable data securely
\item
  Separate names or IDs from survey responses when possible
\item
  Use clear file names and consistent folder structures
\item
  Maintain a codebook or data dictionary
\item
  Document how variables were created or changed
\item
  Limit access to sensitive data
\item
  Back up files in secure locations
\item
  Remove or mask identifying information before sharing data
\item
  Follow institutional review board, institutional, legal, and community
  requirements
\end{itemize}




\end{document}
